\documentclass[conference, a4paper, 10pt]{IEEEtran}

\usepackage{url}
\usepackage{cite}
\usepackage{amsmath,amssymb,amsfonts}
\usepackage{xcolor}

\usepackage{booktabs}
%\usepackage{enumitem}
\usepackage[hidelinks]{hyperref}
%\addbibresource{OurBibliography.bib}

% correct bad hyphenation here
\hyphenation{op-tical net-works semi-conduc-tor}

\begin{document}
\title{Project Title Here}

% author names and e-mail addresses
\author{
\IEEEauthorblockN{Guilherme Ferreira}
\IEEEauthorblockA{up201706719@fe.up.pt}
\and
\IEEEauthorblockN{João Oliveira}
\IEEEauthorblockA{up202205302@fe.up.pt}
\and
\IEEEauthorblockN{Ayrton da Cruz}
\IEEEauthorblockA{up...@fe.up.pt}
\and
\IEEEauthorblockN{Ana Lúcia Ferreira}
\IEEEauthorblockA{up....@fe.up.pt}
}
% generate the title area
\maketitle

% As a general rule, do not put math, special symbols or citations in the abstract
\begin{abstract}
The abstract goes here. Abstracts typically describe the context of the project, the problem that is being solved, the approach at solving it, and identify major results.  This paper-style report is divided in five sections. These are not mandatory, but give you a hint about what you should talk about in the paper. Page limit is 4 two-column pages. If there is a specific reason why you should need more, please let me known. If you want to provide manual-style information e.g. describing how to set up your system do it in another document, put that document online, and cite it here.

For the project proposal, you do not need to include the five sections. You may provide only a small introduction, and perhaps a few links to tools, research articles, or websites that you plan to explore further. 

\end{abstract}

\section{Introduction}
Here you should provide the context of this work and describe the problem you are addressing. You can explain why it is a relevant issue, and describe the approach that you intend to use.

\section{Related work}
You should identify and analyze related work here. List some sources that you have found and that you plan to explore or replicate. Use bibliographic references when possible, and cite them like this~\cite{ArfaouiPrivacy2019}.

Keep your references in a \texttt{.bib} file. You can fill it by copying BibTeX entries directly (e.g., from IEEE Xplore, ACM Digital Library, or even Google Scholar), or you can use a reference manager like Zotero\footnote{Available in \url{https://www.zotero.org/}} for your convenience.

% These sections are commented out because they are not needed for the initial proposal and related work discussion
%\section{Solution}
%Here you should describe the solution.

%\section{Evaluation}
%The evaluation results go here.

%\section{Conclusion}
%The conclusion goes here.


% References section
\bibliographystyle{./IEEEtran}
\bibliography{./IEEEabrv,./OurBibliography}
\end{document}
